


\begin{figure}[htb!]
\begin{center}
\includegraphics[width=0.85\textwidth]{model/figures/clean_chunk_images/clean_chunk_merged_v2.pdf}
\end{center}
\caption{The figure shows how different tasks can be unified by varying the proportion of clean chunks. Each vertical bar represents a latent frame in a chunk, with darker bars indicating higher noise levels and the \emph{white} bars denoting clean frames. The first row illustrates the early inference stage of T2V generation, starting from a single fully noisy chunk and progressing to multiple noisy chunks, before any clean chunk has been produced. The middle row depicts the case of I2V generation, treated as a special case of continuation in which only the first frame of the first chunk is clean. The last row describes a general stage where clean chunks are already available, applicable to video continuation and other scenarios involving prior denoised content.}
\label{fig:clean_chunk}
\end{figure}
